% SHARD-ID: CA-33-GAUSSIAN-STRESS-NUMERICAL-SUITE
% SHARD-TITLE: First Gaussian Stress-Energy Numerical Suite
% SHARD-SUMMARY: Extends the checked Gaussian current layer with finite open-chain T_01 continuity tests across several chains and endpoint signs.
% SHARD-SUMMARY: Adds a one-dimensional nearest-neighbour current-symbol slope residual that detects wrong speed, bond magnitude, and orientation.
% SHARD-SUMMARY: Keeps periodic first moments, finite-chain dGamma(k), higher-dimensional current densities, and T_11 witnesses as open obligations.
% SHARD-KEYWORDS: Gaussian boson, stress-energy numerical suite, open chain, energy current, slope residual, failure witness

\section{First Gaussian Stress-Energy Numerical Suite}
\label{sec:gaussian-stress-numerical-suite}

\subsection{Status, Sources, and Dependencies}

\textbf{Status: checked local-current numerical suite.}  This shard records the
first stress-energy-specific numerical tests.  The checked surface is narrow:
finite open-chain \(T_{01}\) continuity for the CA-29/CA-30 density-current
choice, and a one-dimensional nearest-neighbour current-symbol slope residual.
It is not a finite boost construction, not a \(\mathrm d\Gamma(k)\) theorem, not
a periodic first-moment test, and not a continuum stress-tensor claim.

Dependencies: CA-29--CA-32 and CONVENTIONS.md (j), (k), (l).

% Convention: CONVENTIONS.md:196--239 -- finite periodic Gaussian checks do not
%   define periodic first moments or finite coordinate generators.
% Convention: CONVENTIONS.md:260--269 -- boost-time symbol
%   (1/2) partial_a omega(k)^2 and its sign convention.
% Convention: CONVENTIONS.md:312--347 -- Poincare vector fields translated to
%   self-adjoint commutator signs.
% Convention: CONVENTIONS.md:349--421 -- Gaussian density split, finite open
%   current orientation, endpoint convention, subtraction, and improvement
%   policies.
% Source: report/sections/29_gaussian_boson_real_space_density.tex:74--105 --
%   equal-endpoint split and finite open-boundary policy.
% Source: report/sections/30_gaussian_boson_1d_stress_candidates.tex:83--146 --
%   T_00, T_01, T_10, and T_11 roles and endpoint signs.
% Source: report/sections/31_gaussian_current_symbol_equivalence.tex:55--98 --
%   one-dimensional nearest-neighbour current-symbol equality.
% Source: report/sections/32_gaussian_higher_dimensional_cell_currents.tex:70--132 --
%   higher-dimensional cell-current equations remain proposal-level inputs.
% Checked: src/GaussianBosonCurrents.jl:126--235 -- open-chain currents,
%   closed forms, current-symbol helper, slope residual, and continuity
%   residuals.
% Checked: test/runtests.jl:236--286 -- strengthened open-chain
%   "Gaussian open-chain energy current continuity" tests.
% Checked: test/runtests.jl:324--397 -- current-symbol equality and
%   "Gaussian stress-energy current slope witnesses" tests.

\subsection{Finite Open-Chain Continuity}

The first suite strengthens the A4 local-current witness instead of replacing
it.  For each tested open chain, the Julia layer constructs the CA-29
equal-endpoint site energies \(e_j\), the total Hamiltonian
\[
  H=\sum_{j=1}^N e_j,
\]
and the edge currents
\[
  J_{j+1/2}=i[e_j,e_{j+1}],
  \qquad 1\le j<N .
\]
The residual helper checks
\[
  i[H,e_j]-(J_{j-1/2}-J_{j+1/2})=0
\]
with endpoint currents \(J_{1/2}=J_{N+1/2}=0\).

The test cases now cover several chain sizes and onsite/bond choices:
\[
  (N,a,b)\in
  \{(2,0.75,-0.5),(3,2.0,1.25),(5,7.25,-1.5),(7,0.1,-0.875)\}.
\]
For each case the suite asserts:
\[
\begin{array}{ll}
1.&\sum_j e_j=H,\\
2.&H_{qq}\text{ is the expected tridiagonal matrix with diagonal }a
    \text{ and bond }b,\\
3.&H_{pp}=I,\\
4.&J_{j+1/2}\text{ equals the closed form }
    \frac b2(q_{j+1}p_j-q_jp_{j+1}),\\
5.&J_{j+1/2}\text{ is independent of the onsite coefficient},\\
6.&\text{all continuity residual matrices vanish},\\
7.&i[H,e_1]=-J_{3/2},\quad
   i[H,e_j]=J_{j-1/2}-J_{j+1/2},\quad
   i[H,e_N]=J_{N-1/2}.
\end{array}
\]
The explicit endpoint tests are the important stress-energy addition: boundary
signs are part of the local \(T_{01}\) claim, not discarded finite-size errors.
The suite also checks that a one-site chain cannot produce a bond current and
that an out-of-range open bond index raises an error.

\subsection{Current-Symbol Slope Residual}

The second suite is the smallest symbol-level current-density failure witness
currently justified by CA-31.  For the one-dimensional nearest-neighbour symbol
\[
  \mathcal J_\epsilon(k)=-\epsilon b\sin(\epsilon k),
\]
the small-\(k\) slope is
\[
  \mathcal J_\epsilon'(0)=-\epsilon^2 b.
\]
The new helper returns the speed residual
\[
  \rho_J(b,\epsilon,c)=-\epsilon^2 b-c^2 .
\]
It validates positive spacing, positive target speed, and real bond data before
returning a number.

For the speed-\(c\) nearest-neighbour Klein-Gordon bond
\[
  b=-c^2\epsilon^{-2},
\]
the residual is zero.  The tests assert this for several spacings and speeds.
They also assert failure witnesses:
\[
  b=-\epsilon^{-2}\quad\text{fails when }c\ne1,
\]
wrong bond magnitude gives the expected nonzero residual, and reversing the
bond sign gives
\[
  \rho_J(c^2\epsilon^{-2},\epsilon,c)=-2c^2 .
\]
These are current-density symbol failures for the CA-31 orientation.  They are
not Hessian-only failures and they do not use a periodic coordinate branch.

\subsection{What This Falsifies}

The suite falsifies three classes of local-current mistakes:
\[
\begin{array}{ll}
1.&\text{incorrect open-boundary signs in }
   i[H,e_j]=J_{j-1/2}-J_{j+1/2},\\
2.&\text{wrong orientation in }
   J_{j+1/2}=i[e_j,e_{j+1}],\\
3.&\text{nearest-neighbour current symbols whose low-energy slope does not
   match the selected }c^2.
\end{array}
\]
The onsite-independence checks also catch an accidental insertion of the mass or
onsite term into the nearest-neighbour current.

\subsection{Open Obligations}

The suite deliberately leaves the following undone:
\[
\begin{array}{ll}
1.&\text{finite open-chain wave packets and comparison with }
   \mathrm d\Gamma(k),\\
2.&\text{periodic first moments or periodic coordinate generators},\\
3.&\text{higher-dimensional current-density formulas and incidence-sign tests},\\
4.&\text{next-nearest-neighbour or doubler current-density witnesses},\\
5.&\text{momentum-density }T_{10}\text{ and stress-current }T_{11}\text{
   continuity tests},\\
6.&\text{scaling maps, state choices, domains, and continuum convergence.}
\end{array}
\]
The doubler evidence remains the CA-28 coefficient/minimum-data rejection.  It
should not be described as a current-density failure until a current formula for
that non-nearest-neighbour symbol is derived and tested.
